\uuid{9DYh}
\titre{Intégration par parties}
\auteur{hanani ; mzari}
\organisation{exo7}
\licence{CC BY-NC-SA 4.0}
\datecreate{2020-01-20}
\niveau{L1}
\module{Analyse}
\chapitre{Calcul d'intégrales}
\sousChapitre{Autre}
\difficulte{3}
\theme{Calculs d'intégrales, intégration par parties}
\exo7qcmid{lille-2.integrales.016}

\contenu{
\begin{question}
\qclassification{127.99}
\qtags{motcle=intégration par parties}
On note $\displaystyle I=\int _{0}^{2}x^2\ln (x+1)\, \mathrm{d}x$. Parmi les affirmations suivantes, cocher celles qui sont vraies :
\begin{answers}
\good{$\displaystyle I\leq 4\int _{0}^{2}\ln (x+1)\, \mathrm{d}x$.}
\bad{$\displaystyle I\geq \ln 3\int _{0}^{2}x^2\, \mathrm{d}x$.}
\bad{$\displaystyle I=\frac{8}{3}\ln 3+\frac{1}{3}\int _{0}^{2}\frac{x^3\, \mathrm{d}x}{x+1}$.}
\good{$\displaystyle I=3\ln 3-\frac{8}{9}$.}
\end{answers}
\vskip3mm
\begin{explanations}
Pour tout $x\in ]0,2[$, $x^2\ln (x+1)< 4\ln (x+1)$ et $x^2\ln (x+1)<x^2\ln 3$. Donc $\displaystyle I<4\int _{0}^{2}\ln (x+1)\, \mathrm{d}x$ et $\displaystyle I<\ln 3\int _{0}^{2}x^2\, \mathrm{d}x$. Enfin, une intégration par parties avec $u=\ln (x+1)$ et $v=x^3/3$ donne :
$$I=\frac{8}{3}\ln 3-\frac{1}{3}\int _{0}^{2}\frac{x^3\, \mathrm{d}x}{x+1}=3\ln 3-\frac{8}{9}.$$
\end{explanations}
\end{question}
}
